% ============================================================
% 中間発表要旨テンプレート
%
% 注意：
% - 原則として2ページに収めること。
% - コメント（% で始まる行）は提出前に残っていても問題ない。
% ============================================================

% オプションの説明：
% no-page-number：ページ番号を表示しない（ページ番号を表示する場合はこのオプションを指定しない）
% english：英語で書く場合に指定する．
% dvipdfmx: uplatex でコンパイルする場合に指定する．（一緒に、latexmkrcのpdf_modeを3にすること）
\documentclass[no-page-number]{nanzan-abstract}
\usepackage{float}             % 図や表の位置を強制するためのパッケージ
% 以下はnanzan-abstract.clsで読み込んでいるのでここでは不要。
% amsthm, graphicx, enumitem
% mathtools (amsmathはmathtoolsで読み込まれる)
% 数式フォント関係も設定済み

% ------------------------------------------------------------
% コンパイルエンジンによる設定の切り替え
% ------------------------------------------------------------

\usepackage{iftex}

\ifLuaTeX
  % LuaLaTeX
  \usepackage{hyperref}
\else
  % upLaTeX + dvipdfmx
  \usepackage[dvipdfmx]{hyperref}
  \usepackage{pxjahyper}
\fi

\usepackage{booktabs}
\usepackage{url}
\usepackage{cite}

% ============================================================
% タイトル
% ============================================================

\title{卒業研究の題目を書く}
\subtitle{必要であれば副題を書く} % 副題がない場合はこの行をコメントアウトする．
\author{2026TS000 氏名}
\professor{南山太郎}

\begin{document}

\maketitle


% ============================================================
\section{はじめに}
% ============================================================

% この節では、研究の背景から研究目的までを簡潔に説明します。
%
% 次の流れを意識してください。
%
% 1. 研究対象となるシステムや技術はどのようなものか
% 2. そこにどのような問題があるか
% 3. 本研究では、その問題に対して何を行うのか
% 4. 中間発表時点ではどこまで行ったのか
%
% 単に「○○について研究します」とするのではなく、
% 「なぜその研究が必要なのか」が分かるようにしてください。
%
% 研究目的では、可能な限り、
%
% - 何を実現するのか
% - 何を比較するのか
% - 何を明らかにするのか
%
% まで具体的に記述してください。
%
% また、
%
% - 卒業研究として最終的に目指していること
% - 中間発表までに実際に行ったこと
%
% を区別して書いてください。

ここに本文を書いてください。


% ============================================================
\section{研究背景}
% ============================================================

% この節では、本研究を理解するために必要な技術や
% 既存研究について説明します。
%
% 主に以下について記述してください。
%
% - 研究対象となる技術やシステム
% - 既存の一般的な方法
% - 関連する既存研究
% - 既存方法では十分に扱えていない問題
% - その問題が本研究につながる理由
%
% 技術そのものの説明だけで終わらず、
% 「なぜ既存の方法だけでは不十分なのか」を示してください。
%
% 既存研究や公式ドキュメントについて述べる場合は、
% 必ず参考文献を引用してください。
%
% ------------------------------------------------------------
% 参考文献の引用例
% ------------------------------------------------------------
%
% references.bib に
%
%   @misc{kubernetes, ...}
%
% と登録している場合、
%
%   \cite{kubernetes}
%
% と書いて引用できます。
%
% 複数の文献をまとめて引用する場合は、
%
%   \cite{kubernetes,marchese2022}
%
% のように書きます。
%
% references.bib に登録しただけでは、通常、
% 参考文献一覧には表示されません。
% 本文中で \cite{...} を使用してください。

Kubernetesでは、コンテナ化されたアプリケーションの
配置や管理を自動化できます\cite{kubernetes}。
また、マイクロサービス間の通信を考慮した配置方式も
研究されています\cite{marchese2022}。

% 上の2文は引用方法を示すサンプルです。
% 自分の研究内容に置き換え、不要な文献は削除してください。

ここに自分の研究背景を書いてください。


% ============================================================
\section{研究概要}
% ============================================================

% この節では、卒業研究全体として何を行うのかを説明します。
%
% 次の内容が分かるようにしてください。
%
% - 研究の対象
% - どのような問題を解決するのか
% - どのような方式・システムを用いるのか
% - システムへの入力
% - どのような処理を行うのか
% - どのような出力や結果を得るのか
% - 何を基準として比較・判断するのか
% - 最終的に何を明らかにしたいのか
%
% この節では、
% 「卒業研究として最終的に実現したいもの」の全体像を
% 説明してください。
%
% 実装方法の細かな説明や実験条件は、この節には書かず、
% 第4節、第5節で説明してください。
%
% 可能であれば、システム構成図や処理の流れを示す図を
% 用いてください。
%
% ただし、中間発表時点で未実装の内容については、
%
% 「今後実装する」
% 「最終的には○○する」
%
% などと書き、実装済みの内容と区別してください。

ここに本文を書いてください。


% ---- 図の例 ------------------------------------------------
%
% \begin{figure}[tb]
%   \centering
%   \includegraphics[width=0.8\linewidth]
%                   {figures/system-overview.pdf}
%   \caption{研究システムの概要}
%   \label{fig:system}
% \end{figure}
%
% 本文中では、
%
%   図\ref{fig:system}に研究システムの概要を示します。
%
% のように参照してください。


% ============================================================
\section{システムの実装}
% ============================================================

% この節では、第3節で説明した研究概要のうち、
% 中間発表までに実際に実装した内容について説明します。
%
% この節では「実装」に集中してください。
% 実験条件や実験結果は第5節に記述します。
%
% 必要に応じて以下について説明してください。
%
% - システム構成
% - 使用したハードウェア
% - 使用したソフトウェアやライブラリ
% - 実装した機能
% - 各コンポーネントの役割
% - 処理の流れ
% - 実装上の工夫
%
% 単に、
%
% 「Pythonを使用した」
% 「Dockerを使用した」
%
% などと使用技術を列挙するのではなく、
% それらを研究システムの中でどのように利用したのかを
% 説明してください。
%
% 最終的に実現する予定のシステムの一部しか
% 実装できていない場合も問題ありません。
%
% その場合は、
%
% 「最終方式の前段階として、現在は○○を実装しています」
%
% のように、実装済みの範囲を明確にしてください。
%
% 実装内容が複数に分かれる場合は、必要に応じて
%
%   \subsection{○○の実装}
%
% のように小節に分けてください。

ここに本文を書いてください。


% ============================================================
\section{実験評価}
% ============================================================

% この節では、第4節で説明したシステムや方式を、
% どのように評価したか、およびその結果から何が分かったかを
% 説明します。
%
% 「実装の説明」は第4節で完結させ、
% この節では「評価」に集中してください。
%
% 原則として、
%
%   5.1 実験方法
%   5.2 実験結果と考察
%
% のように分けて記述してください。


% ------------------------------------------------------------
\subsection{実験方法}
% ------------------------------------------------------------

% この小節では、
% 「何を確認するために、どのような実験を行ったのか」
% を説明します。
%
% 以下について必要なものを記述してください。
%
% - 実験の目的
% - 実験環境
% - 使用するデータや対象
% - 比較する方式・条件
% - 変化させるパラメータ
% - 測定する指標
% - 実験回数
%
% 実験条件だけを列挙するのではなく、
%
% 「この実験によって何を確認しようとしているのか」
%
% が分かるようにしてください。
%
% 例えば、
%
% 「提案方式が負荷の変化に追従できるかを確認するため、
%  負荷条件を○○、○○、○○と変化させて比較します」
%
% のように、研究目的と実験条件を結び付けてください。
%
% 実験条件が多い場合は、表を使用すると分かりやすくなります。

ここに本文を書いてください。


% ---- 実験条件の表の例 --------------------------------------
%
% \begin{table}[tb]
%   \centering
%   \caption{実験条件}
%   \label{tab:conditions}
%   \begin{tabular}{lrr}
%     \toprule
%     条件 & 設定A & 設定B \\
%     \midrule
%     条件1 & 10 & 20 \\
%     条件2 & 30 & 40 \\
%     \bottomrule
%   \end{tabular}
% \end{table}


% ------------------------------------------------------------
\subsection{実験結果と考察}
% ------------------------------------------------------------

% この小節では、実験によって得られた結果を示し、
% その結果から何が分かったかを考察します。
%
% 結果は、可能な限り表やグラフを用いて示してください。
%
% そのうえで、
%
% - 条件によってどのような違いがあったか
% - どの方式がどの条件で良かったか
% - 想定していた結果になったか
% - 想定していなかった現象があったか
% - なぜその結果になったと考えられるか
% - 研究目的に対して何が確認できたか
%
% を説明してください。
%
% 単に、
%
% 「○○という結果になりました」
%
% で終わらず、
%
% 「この結果から○○と考えられます」
%
% まで記述してください。
%
% 十分な実験が完了していない場合でも、
% この小節を空欄にはしないでください。
%
% 現時点で行った予備実験や動作確認について、
%
% - 何を試したか
% - どのような結果になったか
% - 何が確認できたか
% - 何がまだ分かっていないか
%
% を記述してください。
%
% また、今回の実験が最終方式ではなく、
% 中間段階のシステムに対する予備評価である場合は、
% そのことを明記してください。

ここに本文を書いてください。


% ============================================================
\section{今後の計画}
% ============================================================

% この節では、中間発表後から卒業論文提出までに
% 行う予定の内容を書いてください。
%
% まず、現在のシステムや評価に残っている問題点・制約を
% 整理すると書きやすくなります。
%
% そのうえで、
%
% - 何を新たに実装するのか
% - 現在の方式のどこを変更するのか
% - どの条件を追加して実験するのか
% - 何を測定するのか
% - どの方式と比較するのか
% - 最終的に何を明らかにするのか
%
% を具体的に書いてください。
%
% 単に、
%
% 「今後さらに実験します」
% 「精度を向上させます」
% 「システムを改善します」
%
% とするだけでは不十分です。
%
% できるだけ、
%
%   現在の問題点
%       ↓
%   今後行うこと
%       ↓
%   それによって何を評価・確認するか
%
% の対応関係が分かるようにしてください。

ここに本文を書いてください。


% ============================================================
\section{おわりに}
% ============================================================

% この節では、要旨全体を簡潔にまとめます。
%
% 基本的には次の3点を書いてください。
%
% 1. 本研究では何を目的としているか
% 2. 中間発表までに何を行い、何が分かったか
% 3. 今後何を行うか
%
% ここでは新しい実験結果や新しい提案を追加しないでください。
%
% 「はじめに」で示した研究目的に対して、
% 現時点でどこまで進んだのかを振り返る形にすると
% 分かりやすくなります。

ここに本文を書いてください。


% ============================================================
% 参考文献
% ============================================================

\bibliographystyle{junsrt}
\bibliography{references} % .bibファイル名。``.bib'' はつけない

\end{document}
